% =====================================================================
%  MarkItDocs — Starter "Carta formal" (licencia MIT)
%  Carta con membrete, destinatario, asunto y firma. EDITA los datos.
% =====================================================================
\documentclass[11pt,a4paper]{article}

\usepackage{iftex}
\ifPDFTeX
  \usepackage[utf8]{inputenc}
  \usepackage[T1]{fontenc}
  \usepackage{lmodern}
\else
  \usepackage{fontspec}
\fi
\usepackage[spanish,es-noshorthands]{babel}
\usepackage[top=2.8cm,bottom=2.8cm,left=2.8cm,right=2.8cm]{geometry}
\usepackage{xcolor}
\usepackage[hidelinks]{hyperref}

\definecolor{membrete}{HTML}{1F4E79}
\setlength{\parskip}{0.7em}
\setlength{\parindent}{0pt}
\pagestyle{empty}

\begin{document}

% ------------------------------ MEMBRETE ------------------------------
\begin{flushright}
  {\large\bfseries\color{membrete} Tu Nombre o Empresa}\\
  Calle y número · Ciudad\\
  correo@ejemplo.com · +00 000 000 000
\end{flushright}

{\color{membrete}\hrule height 1.5pt}
\vspace{1.2em}

\begin{flushright}
  Ciudad, \today
\end{flushright}

% ---------------------------- DESTINATARIO ----------------------------
\textbf{Nombre del destinatario}\\
Cargo\\
Empresa u organización\\
Dirección

\vspace{1.2em}
\textbf{Asunto:} Escribe aquí el asunto de la carta.

\vspace{0.8em}
Estimado/a señor/a:

Primer párrafo: presenta el motivo de la carta de forma directa. Quién eres,
por qué escribes y qué esperas conseguir con esta comunicación.

Segundo párrafo: desarrolla los detalles — antecedentes, datos concretos,
fechas o condiciones relevantes. Mantén un tono profesional y frases cortas.

Tercer párrafo: cierra con la acción esperada (una reunión, una respuesta,
la confirmación de un acuerdo) y un plazo si aplica.

Sin otro particular, quedo a su disposición y le saludo atentamente.

\vspace{3em}
\rule{6cm}{0.4pt}\\
\textbf{Tu Nombre}\\
Cargo o título

\end{document}
